\noindent\textbf{Energy prediction frameworks for LLMs.}
There has been work on predicting the energy consumption of LLMs~\cite{strubell19,Nguyen-hotcarbon,wilkins2024offlineenergyoptimalllmserving,irene,codecarbon}. Many of these works~\cite{strubell19,Nguyen-hotcarbon,wilkins2024offlineenergyoptimalllmserving} rely on utilization, token counts, or software telemetry, which are coarse-grained features. Consequently, these techniques have been shown to incur high prediction errors~\cite{sustainlp}.

IrEne~\cite{irene} is the closest methodological baseline to our work. IrEne takes into account the relationship between different model components when making predictions. The more recent work on energy prediction~\cite{Jesse-dodge-icml} uses CodeCarbon~\cite{codecarbon}, which estimates energy from software telemetry. However, these prediction techniques target single-GPU execution and in some cases rely on coarse aggregate signals. Thus, they substantially underestimate the energy of parallel inference (\S\ref{sec:model_prediction}). This is the gap that \sys{} addresses.

\smallskip
\noindent\textbf{Energy measurement.}
Another area of research focuses on energy measurement~\cite{lacoste2019quantifyingcarbonemissionsmachine, anthony2020carbontrackertrackingpredictingcarbon, codecarbon, energyVis, mauricio, luccioni2022estimatingcarbonfootprintbloom,dnne} rather than prediction. These measurements rely on hardware monitors (external or integrated such as IPMI) or software-based tools such as NVML. Hardware monitors might not always be available and energy measurements require exclusive access to hardware, which is difficult in multi-tenant environments. Software profilers such as NVML only account for GPU power and are often treated as an energy lower bound~\cite{Jesse-dodge-icml}. We empirically confirm in \S\ref{sec:model_prediction} that GPU-only energy is an inaccurate proxy for full-system energy.

\smallskip
\noindent\textbf{Modeling and predicting latency.}
There have been several recent works~\cite{distsim,dpro,srifty,neusight,habitat,daydream} that predict LLM latency instead of energy. However, predicting energy is a fundamentally different problem. GPUs are optimized for performance, so architectural changes such as more cores or higher clocks often translate relatively predictably to latency, whereas energy depends jointly on execution time, power dynamics, communication, and synchronization. In \S\ref{sec:ablation}, we experimentally show that directly retargeting a recent latency-prediction methodology to energy produces substantially higher error than \sys{}. Other works on improving energy efficiency, such as dynamic voltage and frequency scaling (DVFS)~\cite{kakolyris2024sloawaregpufrequencyscaling,pccl}, are orthogonal to our prediction framework.
